
\documentclass{article}
\usepackage{colortbl}
\usepackage{makecell}
\usepackage{multirow}
\usepackage{supertabular}

\begin{document}

\newcounter{utterance}

\centering \large Interaction Transcript for game `mmlu\_pro', experiment `mmlu\_pro\_math', episode 1148 with reflexion{-}llama{-}3.1{-}8b{-}instruct.
\vspace{24pt}

{ \footnotesize  \setcounter{utterance}{1}
\setlength{\tabcolsep}{0pt}
\begin{supertabular}{c@{$\;$}|p{.15\linewidth}@{}p{.15\linewidth}p{.15\linewidth}p{.15\linewidth}p{.15\linewidth}p{.15\linewidth}}
    \# & \multicolumn{2}{c}{Player} && \multicolumn{2}{c}{Game Master} \\
    \hline

    \theutterance \stepcounter{utterance}  
    & & & \multicolumn{4}{p{0.6\linewidth}}{
        \cellcolor[rgb]{0.9,0.9,0.9}{
            \makecell[{{p{\linewidth}}}]{
                \texttt{\tiny{[P1$\langle$GM]}}
                \texttt{You have attempted this same task before and did not succeed. These are the plans you wrote after those attempts.} \\
\texttt{Trial \#1: New plan: } \\
\\ 
\texttt{Upon reflection, I realize that I didn't carefully read the questions and options before answering. I should have taken more time to understand the problem and the possible solutions before making a choice. } \\
\\ 
\texttt{In the future, I will make sure to:} \\
\\ 
\texttt{1. Read each question carefully and understand what is being asked.} \\
\texttt{2. Take a moment to think about the possible solutions and options before making a choice.} \\
\texttt{3. Eliminate obviously incorrect options and focus on the remaining choices.} \\
\texttt{4. Use my knowledge of math concepts and formulas to help me arrive at the correct answer.} \\
\\ 
\texttt{Specifically, I will pay closer attention to the following:} \\
\\ 
\texttt{* In the first question, I should have noticed that the characteristic of the ring 2Z is the smallest positive integer that divides every element of the ring, which is 0 in this case.} \\
\texttt{* In the second question, I should have recognized that the transformation ST {-} TS is the identity map of V onto itself, as it maps each polynomial to itself.} \\
\texttt{* In the third question, I should have noticed that the greatest negative number in the set B is obtained by finding the minimum value of the expression 7m + 12n = 22, which is {-}4.} \\
\texttt{* In the fourth question, I should have used the concept of differential equations to model the concentration of salt in the tank and arrived at the correct answer.} \\
\texttt{* In the fifth question, I should have used the concept of division to find the number of teams needed, which is 6.} \\
\texttt{* In the sixth question, I should have used the concept of combinations to find the number of ways to select 3 students from a class of 20, which is 1140.} \\
\\ 
\texttt{By following these steps, I hope to improve my performance and arrive at the correct answers in the future.} \\
\\ 
\texttt{Trial \#2: New plan: } \\
\\ 
\texttt{Upon reflection, I realize that I didn't follow the specified format for answering the questions, which led to the game master rejecting my responses. I should have answered each question with "The answer is (X)" where X is the correct letter choice.} \\
\\ 
\texttt{In the future, I will make sure to:} \\
\\ 
\texttt{1. Answer each question with the correct format, i.e., "The answer is (X)".} \\
\texttt{2. Carefully read the question and options before making a choice.} \\
\texttt{3. Take a moment to think about the possible solutions and options before making a choice.} \\
\texttt{4. Eliminate obviously incorrect options and focus on the remaining choices.} \\
\texttt{5. Use my knowledge of math concepts and formulas to help me arrive at the correct answer.} \\
\\ 
\texttt{Specifically, I will pay closer attention to the format and make sure to answer each question correctly.} \\
\\ 
\\ 
\texttt{The following are multiple choice questions (with answers) about math. Answer only with "The answer is (X)" where X is the correct letter choice.} \\
\\ 
\texttt{Question:} \\
\texttt{The symmetric group \$S\_n\$ has \$} \\
\texttt{\textbackslash{}factorial\{n\}\$ elements, hence it is not true that \$S\_\{10\}\$ has 10 elements.} \\
\texttt{Find the characteristic of the ring 2Z.} \\
\texttt{Options:} \\
\texttt{A. 0} \\
\texttt{B. 30} \\
\texttt{C. 3} \\
\texttt{D. 10} \\
\texttt{E. 12} \\
\texttt{F. 50} \\
\texttt{G. 2} \\
\texttt{H. 100} \\
\texttt{I. 20} \\
\texttt{J. 5} \\
\texttt{The answer is (A)} \\
\\ 
\texttt{Question:} \\
\texttt{Let V be the set of all real polynomials p(x). Let transformations T, S be defined on V by T:p(x) {-}> xp(x) and S:p(x) {-}> p'(x) = d/dx p(x), and interpret (ST)(p(x)) as S(T(p(x))). Which of the following is true?} \\
\texttt{Options:} \\
\texttt{A. ST + TS is the identity map of V onto itself.} \\
\texttt{B. TS = 0} \\
\texttt{C. ST = 1} \\
\texttt{D. ST {-} TS = 0} \\
\texttt{E. ST = T} \\
\texttt{F. ST = 0} \\
\texttt{G. ST = TS} \\
\texttt{H. ST {-} TS is the identity map of V onto itself.} \\
\texttt{I. TS = T} \\
\texttt{J. ST = S} \\
\texttt{The answer is (H)} \\
\\ 
\texttt{Question:} \\
\texttt{Let A be the set of all ordered pairs of integers (m, n) such that 7m + 12n = 22. What is the greatest negative number in the set B = \{m + n : (m, n) \textbackslash{}in A\}?} \\
\texttt{Options:} \\
\texttt{A. {-}5} \\
\texttt{B. 0} \\
\texttt{C. {-}3} \\
\texttt{D. {-}7} \\
\texttt{E. {-}4} \\
\texttt{F. {-}6} \\
\texttt{G. {-}1} \\
\texttt{H. {-}2} \\
\texttt{I. {-}9} \\
\texttt{J. N/A} \\
\texttt{The answer is (E)} \\
\\ 
\texttt{Question:} \\
\texttt{A tank initially contains a salt solution of 3 grams of salt dissolved in 100 liters of water. A salt solution containing 0.02 grams of salt per liter of water is sprayed into the tank at a rate of 4 liters per minute. The sprayed solution is continually mixed with the salt solution in the tank, and the mixture flows out of the tank at a rate of 4 liters per minute. If the mixing is instantaneous, how many grams of salt are in the tank after 100 minutes have elapsed?} \\
\texttt{Options:} \\
\texttt{A. 3 + e\^{}{-}2} \\
\texttt{B. 2 {-} e\^{}{-}4} \\
\texttt{C. 2 {-} e\^{}{-}2} \\
\texttt{D. 3 + e\^{}{-}4} \\
\texttt{E. 2 + e\^{}{-}3} \\
\texttt{F. 2 {-} e\^{}{-}3} \\
\texttt{G. 3 {-} e\^{}{-}2} \\
\texttt{H. 2 + e\^{}{-}2} \\
\texttt{I. 2 + e\^{}{-}4} \\
\texttt{J. 2} \\
\texttt{The answer is (I)} \\
\\ 
\texttt{Question:} \\
\texttt{A total of 30 players will play basketball at a park. There will be exactly 5 players on each team. Which statement correctly explains how to find the number of teams needed?} \\
\texttt{Options:} \\
\texttt{A. Multiply 5 by 5 to find 25 teams.} \\
\texttt{B. Divide 30 by 5 to find 6 teams.} \\
\texttt{C. Add 5 to 30 to find 35 teams.} \\
\texttt{D. Subtract 30 from 5 to find {-}25 teams.} \\
\texttt{E. Divide 5 by 30 to find 0.1667 teams.} \\
\texttt{F. Add 5 to 30 then divide by 2 to find 17.5 teams.} \\
\texttt{G. N/A} \\
\texttt{H. N/A} \\
\texttt{I. N/A} \\
\texttt{J. N/A} \\
\texttt{The answer is (B)} \\
\\ 
\texttt{Question:} \\
\texttt{In how many ways can 3 students be selected from a class of 20 to form a study group?} \\
\texttt{Options:} \\
\texttt{A. 560} \\
\texttt{B. 1140} \\
\texttt{C. 1000} \\
\texttt{D. 680} \\
\texttt{E. 760} \\
\texttt{F. 210} \\
\texttt{G. 1500} \\
\texttt{H. 190} \\
\texttt{I. 850} \\
\texttt{J. 1200} \\
            }
        }
    }
     \\ \\

    \theutterance \stepcounter{utterance}  
    & \multicolumn{4}{p{0.6\linewidth}}{
        \cellcolor[rgb]{0.561,0.737,0.561}{
            \makecell[{{p{\linewidth}}}]{
                \texttt{\tiny{[P1$\rangle$GM]}}
                \texttt{The answer is (B)} \\
            }
        }
    }
    & & \\ \\

    \theutterance \stepcounter{utterance}  
    & & & \multicolumn{2}{p{0.3\linewidth}}{
        \cellcolor[rgb]{0.9,0.9,0.9}{
            \makecell[{{p{\linewidth}}}]{
                \texttt{\tiny{[GM$|$GM]}}
                \texttt{B} \\
            }
        }
    }
    & & \\ \\

    \theutterance \stepcounter{utterance}  
    & & & \multicolumn{2}{p{0.3\linewidth}}{
        \cellcolor[rgb]{0.9,0.9,0.9}{
            \makecell[{{p{\linewidth}}}]{
                \texttt{\tiny{[GM$|$GM]}}
                \texttt{B} \\
            }
        }
    }
    & & \\ \\

    \theutterance \stepcounter{utterance}  
    & & & \multicolumn{2}{p{0.3\linewidth}}{
        \cellcolor[rgb]{0.9,0.9,0.9}{
            \makecell[{{p{\linewidth}}}]{
                \texttt{\tiny{[GM$|$GM]}}
                \texttt{game\_result = WIN} \\
            }
        }
    }
    & & \\ \\

\end{supertabular}
}

\end{document}
